\documentclass[12pt,a4paper,UTF8,fontset=windows]{ctexart}

\usepackage[a4paper,margin=2.4cm,bottom=2.6cm]{geometry}
\usepackage{array}
\usepackage{amssymb}
\usepackage{booktabs}
\usepackage{tabularx}
\usepackage{longtable}
\usepackage{enumitem}
\usepackage{xcolor}
\usepackage{hyperref}

\hypersetup{
  colorlinks=true,
  linkcolor=black,
  filecolor=black,
  urlcolor=blue,
  pdftitle={铁程 tripai 设计文档},
  pdfauthor={铁程 tripai 团队}
}

\setlist{nosep}
\renewcommand{\arraystretch}{1.32}

\newcommand{\sectionline}{\par\vspace{0.2em}\noindent\rule{\linewidth}{0.4pt}\par}
\newcommand{\cnsection}[1]{%
  \clearpage
  \begin{center}
    {\Large\bfseries #1}
  \end{center}
  \sectionline
  \addcontentsline{toc}{section}{#1}
}
% 短节连排（不强制换页）：用于内容较少的末尾章节
\newcommand{\cnsectionflow}[1]{%
  \par\vspace{1.8em}
  \begin{center}
    {\Large\bfseries #1}
  \end{center}
  \sectionline
  \addcontentsline{toc}{section}{#1}
}
\newcommand{\subsectioncn}[1]{\par\vspace{0.6em}\noindent{\large\bfseries #1}\par\vspace{0.25em}}
\newcommand{\minisubsection}[1]{\par\vspace{0.4em}\noindent{\bfseries #1}\par\vspace{0.15em}}

\begin{document}
\pagestyle{plain}   % 页码在页脚居中；清除任何页眉标记

\begin{titlepage}
  \centering
  {\Large 2026年华北五省（市、自治区）及港澳台大学生\par}
  \vspace{0.6em}
  {\huge\bfseries 计算机应用大赛\quad 移动互联网应用创新\par}
  \vspace{0.5em}
  {\Large 移动互联网应用程序开发\par}
  \vspace{5.5em}
  {\large 【项目名称】\par}
  \vspace{0.3em}
  {\LARGE\bfseries tripai—原生智能火车出行规划agent\par}
  \vspace{0.8em}
  {\large （Web 应用作品）\par}
  \vspace{6.5em}
  \begin{flushleft}
    \hspace*{4em}
    \begin{tabular}{@{}r@{\ }l@{}}
      所在赛区： & \rule[-0.2em]{3cm}{0.4pt}北京赛区 \rule[-0.2em]{3cm}{0.4pt} \\
      所在学校： & \rule[-0.2em]{2cm}{0.4pt}中国人民大学 \rule[-0.2em]{2cm}{0.4pt} \\
      团队名称： & \rule[-0.2em]{2cm}{0.4pt}人大T0队 \rule[-0.2em]{2cm}{0.4pt} \\
      团队成员： & \rule[-0.2em]{1cm}{0.4pt}温智皓、马霄宇、张泽昌、宋诗琳、朱明哲 \rule[-0.2em]{1cm}{0.4pt} \\
      指导教师： & \rule[-0.2em]{3cm}{0.4pt}张超 \rule[-0.2em]{3cm}{0.4pt} \\
      提交日期： & \rule[-0.2em]{2cm}{0.4pt}2026年10月11日 \rule[-0.2em]{2cm}{0.4pt} \\
    \end{tabular}
  \end{flushleft}
\end{titlepage}

\begin{center}
  {\large\bfseries 2026年华北五省（市、自治区）及港澳台大学生计算机应用大赛}\\[0.35em]
  {\Large\bfseries 参赛作品知识产权声明}
\end{center}

\vspace{0.7em}

\noindent\textbf{参赛作品名称：铁程 tripai——对话式智能火车出行规划}（下称该作品）。

\medskip
\noindent 本小组全体成员，就该作品声明如下：

\begin{enumerate}[label=\arabic*、,leftmargin=2em]
  \item 该作品（含提交参赛的Web应用、相关文档和介绍视频）是本小组成员的原创研究成果，不存在侵犯任何他人知识产权、涉及泄密或违反中华人民共和国法律法规的情形，且未在规定时间前公开发布。
  \item 我们确认：该作品知识产权归本小组成员所有，且签署本声明的人员能够代表全体小组成员的意愿。
  \item 我们承诺：将仅由本小组成员参加大赛答辩。
  \item 大赛举办方为宣传大赛、推广参赛作品，以及为未来各届大赛的参赛选手提供参考等非盈利目的，可能需要以下列形式使用参赛作品，就此我们确认：
  \begin{enumerate}[label=（\arabic*）,leftmargin=3em]
    \item $\square$同意\quad $\square$不同意：将参赛作品的相关文档结集出版（含公开发行）；
    \item $\square$同意\quad $\square$不同意：在相关网站（包括大赛官网及含优酷等第三方视频平台网站）上传和播放介绍视频或幻灯片；
    \item $\square$同意\quad $\square$不同意：在大赛官方网站提供APP应用免费下载；
    \item 作出上述许可，同样不违反本声明第1款的要求。
  \end{enumerate}
\end{enumerate}

\vspace{1.6em}
\noindent 作品小组全体成员：\rule[-0.2em]{7cm}{0.4pt}

\vspace{0.7em}
\noindent 日期：2026年\underline{\hspace{1.1em}}月\underline{\hspace{1.1em}}日

\newpage
\tableofcontents
\markboth{}{}   % 清除“目录”页标记，避免残留到后续页面

\cnsection{一、作品概述}

“铁程 tripai”是一个面向个人出行者的对话式智能火车出行规划系统。传统查询工具通常要求用户逐项填写车站、日期、坐席和换乘条件，遇到直达无票时，用户还要自行拆分区间并逐车核实余票。且市面主流购票并未融合agent，遇到对自身需求模糊的用户有很大的局限性。本作品将这一过程交给AI Agent：用户用一句自然语言描述“何时、几人、从哪里到哪去、坐什么席别、偏好什么时间”等等自然语言，系统自动完成车站消歧、本地线路查询、12306实时余票与票价核实、补票变体分析和正式报告生成。

作品的核心价值是把“本地全路网静态推理”和“联网实时票额确认”分成两层。本地层保存车站、车次、经停和历史票价，可在离线状态下快速枚举方案；联网层只对推荐候选调用12306接口，减少无效请求并保证最终建议可购买、可执行。系统同时保留手动查票模式，满足需求明确用户的精确筛选和快速直查要求。并符合12306使用规范。

\cnsection{二、市场现状、可行性分析与目标群体}

\subsectioncn{（1）市场现状分析}
当前火车出行查询市场呈“一超多强”格局：12306官方渠道权威可靠，但以直达票查询为主，中转推荐能力有限；携程、飞猪、智行等第三方平台虽提供中转方案，却以车票代销与附加服务推销为导向，用户仍需在多个页面间手动筛选比对。行业共性痛点集中在三点：一是余票、时刻、攻略信息分散在多个平台，来回切换成本高；二是查询结果“能看不能判”，是否真正可购、能否衔接成功需用户自行确认；三是购票类应用功能日趋复杂，老年用户面对层层表单与营销入口，数字门槛明显。

本作品正是针对上述缺口：以AI Agent把“理解需求—本地比案—联网核实—给出可购建议”的完整链路自动化，用一次对话替代多平台往返操作。

\subsectioncn{（2）可行性分析}

\minisubsection{技术可行性}
后端采用Python 3和FastAPI，数据处理使用SQLite，前端使用原生HTML/CSS/JavaScript，服务端事件（SSE）承担流式输出。12306公开查询接口可核实余票和票价，OpenAI兼容接口可接入不同大模型。当前本地库包含3404个车站、16076趟车次、162943条经停记录和约279万条历史票价记录，具备支撑全路网方案枚举的数据基础。

\minisubsection{操作可行性}
用户无需注册登录。首次使用只需配置API-KEY、模型名称和接口地址；日常使用可双击启动脚本，浏览器自动打开本地服务。AI模式适合自然语言提问，手动模式提供表单筛选，两类用户均能低门槛上手。

\minisubsection{经济可行性}
系统主体可运行在个人电脑或一台低配置服务器上，SQLite文件部署简单；模型服务由用户选择已有OpenAI兼容账号承担，12306核实接口按需调用，不引入额外数据库许可成本。

\subsectioncn{（3）目标群体}
本作品以AI的便捷性与适老性为首要服务目标。对不熟悉智能手机应用的老年旅客而言，主流购票应用的站内搜索、席位图表、候补规则学习成本高，而本系统只需一句口语化描述——例如“明天上午带老伴从北京去天津看孩子，要靠窗的座”——即可自动完成车站消歧、方案枚举与余票核实，并以大字号、无营销干扰的界面直接给出可执行建议，全程免注册、免下载、零表单操作；想到哪说到哪，AI会主动追问补全。

核心用户还包括：需要在直达无票时寻找中转方案的学生、上班族、商务出行者和跨城通勤人群；需要按时间窗、坐席、车型精确筛选并快速直查的高频用户（手动查票模式）；以及需要保留规划过程、比较多个候选方案的团队出行组织者。

\cnsection{三、作品功能与原型设计}

\subsectioncn{（1）功能概述}
系统面向两类角色：出行旅客（核心服务对象）与数据维护者（采集与运维）。功能需求编号F1～F7如下：

\begin{longtable}{|p{0.9cm}|p{2.6cm}|p{9.7cm}|p{1.3cm}|}
\hline
\textbf{编号} & \textbf{功能} & \textbf{说明} & \textbf{优先级} \\ \hline
F1 & AI对话规划 & 用户用自然语言提出需求；系统展示思考流和工具调用轨迹，流式生成正式规划报告与方案总表，支持连续追问和多轮调整。 & 高 \\ \hline
F2 & 本地方案枚举 & 基于本地车站、车次、经停数据枚举直达、一次换乘和两次换乘方案，执行换乘间隔、当日到达、时间窗、车型等硬约束。 & 高 \\ \hline
F3 & 实时余票核实 & 对推荐方案调用12306接口确认余票和票价，识别不可购区间；无票时继续分析前后多站、买短补长等变体。 & 高 \\ \hline
F4 & 对话详情 & 全屏平铺会话中的全部结构化方案，包括无票方案；支持逐条或批量联网核实，便于用户比较全量候选。 & 中 \\ \hline
F5 & 手动查票 & 不经AI直接连接查票引擎，支持站名模糊匹配、日期、坐席、换乘次数、时间范围、车型、排序和返回数量筛选。 & 高 \\ \hline
F6 & 会话持久化 & 每个会话保存为JSON文件，包含消息、思考链、工具轨迹和方案数据；重启后可恢复，历史会话可完整回放。 & 中 \\ \hline
F7 & 数据采集维护 & 提供车站、车次经停、票价采集和异常清理脚本，支持断点续爬、并发限速、单趟补采和状态统计。 & 中 \\ \hline
\end{longtable}

\minisubsection{非功能性需求}
性能：本地方案枚举秒级返回，联网核实按购买区间合并请求以控制耗时；可用性：AI与手动双模式覆盖不同熟练度用户，长任务全程有进度反馈且可随时中止；可靠性：采集支持断点续爬，接口异常返回可读错误，历史参考价与实时票价严格区分标注；安全性：API-KEY仅保存在本机.env，不代用户下单、不存储任何支付信息。

\subsectioncn{（2）原型设计}
主界面采用“模式选择—具体视图”的两级结构。模式选择页突出AI对话与手动查票两条路径；AI视图左侧为会话列表，中部为思考流、正式回答和方案总表；对话详情页按卡片平铺候选方案，并保留联网核实按钮；手动查票页左侧为筛选表单，右侧为结果列表。整体使用人大红与米白纸感配色，突出正式、稳定和易读的出行助手形象。

交互设计遵循三项原则：第一，同时支持手动后台调用和AI agent辅助；第二，本地枚举结果与联网核实状态分离展示，实现轻量化联网+准确性核查的两步走策略；第三，应用harness工程，实现长任务提供查询中、检查中、思考秒数等反馈。

\cnsection{四、作品实现、难点及特色分析}

\subsectioncn{（1）作品实现及难点}

\minisubsection{1. 总体架构}
系统分为采集层、数据层、规划引擎、AI Agent层、服务接口层和前端表现层。\texttt{collect}包负责静态数据采集；\texttt{data/railway.db}保存车站、车次和经停，\texttt{data/prices.db}保存历史票价；\texttt{app/planner.py}执行纯本地规划；\texttt{app/tools.py}定义Agent可调用工具；\texttt{app/online.py}封装12306实时核实；\texttt{app/server.py}提供REST和SSE接口；\texttt{static}目录实现响应式页面。

\minisubsection{2. 开发周期设计过程}

\begin{longtable}{|p{2.5cm}|p{12.3cm}|}
\hline
\textbf{周期} & \textbf{主要工作与产出} \\ \hline
需求分析 & 梳理直达查询、跨天列车、多次换乘、无票补救、实时票额、多轮对话和会话回放等场景，确定“AI规划”与“手动直查”双模式。 \\ \hline
数据设计 & 设计车站表、车次表、经停表、站车索引和历史票价表；为站对查询、车次编号匹配、改号车别名和换乘枚举建立索引。 \\ \hline
架构设计 & 确定本地枚举与联网核实分离、SQLite轻量部署、SSE流式响应、JSON会话落盘和采集脚本独立运行的技术路线。 \\ \hline
详细设计 & 定义换乘图搜索约束、评分公式、工具协议、方案数据结构、票额状态、异常提示和前端卡片层级。 \\ \hline
编码实现 & 实现采集器、规划引擎、Agent工具循环、12306核实、会话管理、前端渲染与批量核实接口。 \\ \hline
测试与修正 & 针对大城市/小站、直达无票、一次换乘、两次换乘、跨零点、非每日车、改号车和接口反爬场景做功能验证。 \\ \hline
部署运维 & 提供启动脚本、采集脚本、日志目录、断点续爬和状态统计；支持按字头小规模试采和全量重建。 \\ \hline
\end{longtable}

\minisubsection{3. 关键实现}
本地后端方案筛选作为独立模块，同时服务于AI agent和手动直查。此模块和联网查询、基础信息了解等tools，共同构成了agent的工具集，轻量化调用结果注入实现了良好的上下文长度管理，结合文本渲染、格式化输出等措施，实现了易于使用的直接基于购票系统的AI agent。 

综合评分由用时、出发时间、换乘代价和等待时间组成，直达不计算等待惩罚。AI工具从约$N \times 20$个去重候选中选择前$N$项供模型分析，同时把前$2N$项送入前端详情页，兼顾回答聚焦与全量可比。

实时核实层将12306返回的坐席、余票、票价映射为统一状态，按\texttt{train\_no}兜底改号车，并对非每日车尝试多个日期偏移。多区间请求会合并查询，减少接口压力。前端使用\texttt{marked}渲染Markdown，并用\texttt{DOMPurify}净化输出。

\minisubsection{4. 主要难点}
\begin{itemize}
  \item 全路网换乘组合规模大：通过层间截断、累计历时限制、城市过滤和配额控制保证响应速度。
  \item 12306数据语义复杂：对电报码、显示车次、内部\texttt{train\_no}、多日期开行和部分区间票额建立统一解析规则。
  \item AI结果需要稳定：模型只负责理解需求、组织工具调用和解释结果，路线计算与票额判定由确定性代码完成。
  \item 长任务体验：SSE分段输出思考流与答案，工具轨迹可展开，用户能观察进度并主动停止。
\end{itemize}

\subsectioncn{（2）特色分析}
\begin{itemize}
  \item \textbf{原生Agent闭环：}不是简单让模型推荐车次，而是模型连续调用本地规划工具、在线核实工具和补票变体工具，结论由数据和代码支撑。
  \item \textbf{“全量＋推荐”双视图：}正式报告只展示可购或明确标记的方案，对话详情保留全部候选，避免用户丢失被AI摘要过滤的信息。
  \item \textbf{补票变体深挖：}直达无票时自动分析前后站、买短乘长和分段组合，并明确区分“可购”“无票/受限”“未核实”。
  \item \textbf{轻量可维护：}SQLite文件库、脚本化采集、断点续爬和日志留存使系统无需复杂数据库服务也能长期维护。
\end{itemize}

\subsectioncn{（3）测试账号与配置说明}

\begin{center}
\fbox{%
\begin{minipage}{0.92\linewidth}
  \textbf{本软件不设置注册用户名和密码，无需测试账号。}首次进入需在本机配置用户自己的API-KEY、模型名称和OpenAI兼容接口地址；该配置保存在项目目录的\texttt{.env}文件中，不随作品源码提交。若评审环境没有模型服务，可先使用手动查票模式体验本地方案枚举，联网核实功能另需可访问12306的网络。
\end{minipage}%
}
\end{center}

\subsectioncn{（4）操作步骤及说明}

\minisubsection{步骤一：启动系统}
首次安装依赖后，双击项目根目录的\texttt{启动铁程.bat}；或在命令行进入项目目录并执行\texttt{.venv\textbackslash Scripts\textbackslash python.exe -m app.server}。浏览器访问\texttt{http://127.0.0.1:8600}。

\minisubsection{步骤二：配置模型}
在配置门页填写API-KEY、模型名称和接口地址，点击“验证并进入”。配置只保存在本机\texttt{.env}；后续可在顶栏“设置”修改。

\minisubsection{步骤三：AI对话规划}
选择“AI对话模式”，在输入框描述需求，例如“下周六3个人从北京去上海，尽量上午出发，二等座”。发送后系统显示思考流、工具调用和正式报告；可在方案总表中查看评分、时刻、票价、余票状态，并继续追问或改期。

\minisubsection{步骤四：查看对话详情}
点击顶栏“对话详情”查看本次会话全部结构化方案。输入前$N$条数量后点击“批量核实”，或对单条方案点击“联网核实”。核实完成后，状态会更新为可购、无票/受限或其他提示。

\minisubsection{步骤五：手动查票}
点击印章图标返回模式选择，进入“手动查票模式”。填写或选择出发站、到达站、出行日期、坐席、换乘次数、时间范围、车型和排序；点击“查询方案”查看本地候选。结果右侧提供“联网核实”，可获得12306实时余票和票价。

\minisubsection{步骤六：会话与数据维护}
AI视图左侧可新建会话；右键会话可重命名或删除。切换历史会话可回放思考链和方案。数据维护可在项目根目录执行\texttt{python -m collect.collector --stats}、\texttt{python -m collect.price\_collector --stats}查看状态，使用\texttt{--resume}续爬；全量采集建议按README中的车站、车次、票价顺序执行。

\subsectioncn{（5）测试设计与质量保障}
测试覆盖五类路径：车站消歧测试（北京、北京南、上海虹桥等同城多站）、约束测试（预售期、跨零点、换乘间隔）、枚举测试（直达、一转、两转及无解场景）、接口测试（配置校验、会话创建、SSE中断、批量核实）和数据测试（改号车别名、非每日车、票价缺失回落）。前端重点验证Markdown净化、长列表滚动、状态标签颜色和会话回放一致性。所有网络失败均返回可读错误，不把历史参考价误标为实时票价。

\newpage

\cnsectionflow{五、团队介绍和人员分工}

\begin{longtable}{|p{3cm}|p{14cm}|}
\hline
\textbf{成员} & \textbf{分工} \\ \hline
张泽昌 & 需求分析、总体架构、AI Agent设计、规划引擎联调、文档统稿。 \\ \hline
马霄宇 & 云端部署、技术支持、爬虫分析、硬件设备支持。 \\ \hline
温智皓 & 数据模型设计、采集器开发、断点续爬与数据清理、索引优化。 \\ \hline
朱明哲 & 12306接口封装、票额解析、多日期兜底、批量核实与异常处理。 \\ \hline
宋诗琳 & 前端界面、SSE渲染、方案卡片、对话详情、交互状态与可用性测试。 \\ \hline
指导教师 & 总体方案指导、技术路线评审与作品完善建议。 \\ \hline
\end{longtable}


\cnsectionflow{六、其他}

\minisubsection{原创性说明与第三方开源组件}
本作品全部业务代码——包括AI Agent对话循环、本地方案规划引擎、12306实时核实层、数据采集器、Web服务接口与前端交互逻辑——均为团队成员原创编写，未复制、未修改或拼装任何第三方开源项目的源代码。

开发过程中仅以标准依赖包形式引用以下公开发布的开源软件（均通过pip或CDN直接获取，未改动其源码），按要求标明如下：

\begin{longtable}{|p{3.4cm}|p{3.2cm}|p{2.8cm}|p{5.8cm}|}
\hline
\textbf{组件} & \textbf{许可证} & \textbf{引入方式} & \textbf{用途} \\ \hline
FastAPI & BSD-3-Clause & pip 依赖 & Web 服务框架与接口定义 \\ \hline
Uvicorn & BSD-3-Clause & pip 依赖 & ASGI 服务器 \\ \hline
Requests & Apache-2.0 & pip 依赖 & HTTP 客户端（12306 接口调用） \\ \hline
Pydantic & MIT & pip 依赖 & 数据模型与参数校验 \\ \hline
SQLite & Public Domain & Python 标准库 sqlite3 & 本地数据存储 \\ \hline
marked.js & MIT & CDN 引入 & 前端 Markdown 渲染 \\ \hline
DOMPurify & Apache-2.0 / MPL-2.0 & CDN 引入 & 前端输出净化（防 XSS） \\ \hline
\end{longtable}

此外，系统调用的12306公开查询接口与OpenAI兼容的大模型接口均为网络服务协议，不涉及开源代码引用；本文档使用LaTeX（TeX Live 发行版）排版。

本作品使用的数据来自公开查询接口采集与整理，仅用于出行研究和个人规划；实时余票和票价仍以12306及车站官方销售结果为准。系统不代用户下单，不存储支付信息，不绕过身份认证或风控规则。采集器内置等待和重试机制，使用时应控制频率，遵守相关网站服务条款。

若部署到公网，建议增加反向代理、HTTPS、访问控制和模型密钥隔离；当前默认版本面向个人本机使用，因此未内置多用户账号体系。

\cnsectionflow{七、致谢}
感谢指导教师在需求收敛、系统架构上的指导；感谢团队成员在数据采集、引擎调试和界面打磨中的协作；也感谢开源社区和学校提供的实验环境。团队后续将继续优化换乘体验、数据更新效率和模型解释能力，使作品更贴近真实出行决策。

\end{document}

